\documentclass[11pt]{amsart}
\usepackage{amssymb}
\usepackage[margin=1.15in]{geometry}
\usepackage[hidelinks]{hyperref}
\newtheorem{theorem}{Theorem}
\newtheorem{lemma}[theorem]{Lemma}
\newtheorem{corollary}[theorem]{Corollary}
\theoremstyle{remark}
\newtheorem{remark}[theorem]{Remark}
\newcommand{\code}[1]{\texttt{#1}}
\newcommand{\Gr}{\mathrm{Gr}}
\newcommand{\adj}{\mathrm{adj}}
\title{Wiebe's theorem relative to an ideal, and its Unit formalisation}
\date{October 25, 2026}
\begin{document}
\maketitle

This note presents a variation of the proof of Wiebe's theorem in
T.~Coquand and C.~T\^ete, \emph{An elementary proof of Wiebe's Theorem}
(J.~Algebra 499, 2018; \code{ctete.pdf}, Theorem 11 and Corollary 12), and
its checked formalisation in Unit (\code{sag/wiebe}).  The argument is
the one of the paper.  Three points are made explicit: everything is
stated relative to a finitely generated ideal $J$ instead of in quotient
rings; the generic change of the sequences is a finite chain of
elementary operations, each of which visibly keeps the determinant; and
the rings obtained by adding indeterminates enter only through two
properties of the inclusion maps.

\section{Statement}

Let $R$ be a commutative ring and $J$ a finitely generated ideal.  For a
list $c=(c_1,\dots,c_n)$ write $(J,c)$ for the ideal generated by $J$
and the $c_i$.  An element $r$ is \emph{regular mod $J$} if $rx\in J$
implies $x\in J$; a list is regular mod $J$ if $c_ix\in J$ for all $i$
implies $x\in J$; $c_1,\dots,c_k$ is a regular sequence mod $J$ if each
$c_i$ is regular mod $(J,c_1,\dots,c_{i-1})$.  The true grade
$\Gr(I)\geq k$ mod $J$ of a finitely generated ideal $I$ (ctete,
Definition 5, relative to $J$) means that, merely, after adding finitely
many indeterminates, $I$ contains a regular sequence of length $k$
mod $J$.  Theorem 9 of ctete holds in this form: if $\Gr(I)\geq k+1$
mod $J$ and $r\in I$ is regular mod $J$, then $\Gr(I)\geq k$ mod $(J,r)$.

\begin{theorem}\label{thm:wiebe}
Let $a,b\in R^n$ and let $V$ be an $n\times n$ matrix with
$b_j\equiv\sum_i a_iV_{ij}$ mod $J$ for all $j$, and $\Delta=\det V$.
\begin{enumerate}
\item For every $i$, $\Delta a_i\in(J,b)$.  Hence $x\in(J,b,a)$ implies
  $x\Delta\in(J,b)$, and $x\in(J,b,\Delta)$ implies $xa_i\in(J,b)$ for
  all $i$.
\item If $\Gr(b)\geq n$ mod $J$, the converses hold: $x\Delta\in(J,b)$
  implies $x\in(J,b,a)$, and $xa_i\in(J,b)$ for all $i$ implies
  $x\in(J,b,\Delta)$.
\end{enumerate}
So, modulo $(J,b)$, $\Delta^\perp=(a)$ and $(a)^\perp=(\Delta)$.
\end{theorem}

With $J=0$ this is ctete Theorem 11.  Since a regular sequence $b$
mod $J$ of length $n$ gives $\Gr(b)\geq n$ mod $J$ directly, it contains
Corollary 12.

Part 1 is Cramer's rule: writing $b\equiv aV$ for row vectors,
$b\,\adj(V)\equiv aV\adj(V)=\Delta a$ mod $J$, so
$\Delta a_k\equiv\sum_j b_j\adj(V)_{jk}\in(b)$.  The rest of this note
proves part 2, by induction on $n$, for all rings $R$ and ideals $J$ at
once.  For $n=0$, $\Delta=1$ and both statements are trivial.

\section{Three invariance properties}

\paragraph{Elementary operations.}  For an index $r$ and coefficients
$w$ with $w_r=0$, the elementary operation $c\mapsto c'$ adds
$\sum_lw_lc_l$ to $c_r$ and leaves the other entries.  It is undone by
the operation with $-w$, since $\sum_l w_lc'_l=\sum_lw_lc_l$ (because
$w_r=0$), so $(J,c')=(J,c)$.  Its matrix $E$ has determinant $1$: up to
transposition it is the identity with one column replaced by
$e_r+w$, and $\det I[r\leftarrow v]=v_r$ because
$v\mapsto\det I[r\leftarrow v]$ is linear and takes the value
$\delta_{kr}$ at $e_k$.
\begin{itemize}
\item An operation on $b$ replaces $V$ by $VE$ (in the row-vector
  convention), so $b'\equiv aV'$ mod $J$ with $\det V'=\Delta$.
\item An operation on $a$ gives $b\equiv a'E^{-1}V$, with
  $\det(E^{-1}V)=\Delta$, and $(a')=(a)$.
\end{itemize}
The conclusions of part 2 only depend on the ideals $(J,a)$, $(J,b)$ and
on $\Delta$: if $(J,b')=(J,b)$ and every $a'_i$ is in $(J,a)$, the
conclusions for $(a',b',\Delta)$ give those for $(a,b,\Delta)$.  So part 2
for $(a,b,V)$ follows from part 2 for any data reached from it by
elementary operations.

\paragraph{Good ring maps.}  Call $h:R\to T$ \emph{good} if it keeps
regularity mod lists ($r$ regular mod $L$ implies $h(r)$ regular mod
$h(L)$) and has a retraction $\rho:T\to R$ (a ring map with
$\rho h=\mathrm{id}$).  The inclusion $R\to R[X]$ is good (McCoy for the
first property, evaluation at $0$ for the second), and composites of
good maps are good.  Along any ring map, $b\equiv aV$ mod $J$ becomes
$h(b)\equiv h(a)h(V)$ mod $h(J)$ and $\det h(V)=h(\Delta)$; along a good
map, the conclusions of part 2 for $h(J),h(a),h(b),h(\Delta)$ give those
for $J,a,b,\Delta$, since $\rho$ brings memberships of images back.

\paragraph{Changing the ideal.}  If $\Gr(I)\geq k$ mod $J$ and every
element of $I$ is in $(J,I')$, then $\Gr(I')\geq k$ mod $J$: each term of
a regular sequence in $I$ is $t+j$ with $t\in(I')$, $j\in J$, and a
regular sequence mod $J$ stays one when its terms are changed mod $J$.

\section{The generic position}

\begin{lemma}\label{lem:generic}
Let $\Gr(b_1,\dots,b_n)\geq n$ mod $J$.  Then, merely, along a good map
$h:R\to T$, the list $h(b)$ is changed by elementary operations, not
touching $b_n$, into $b'$ such that $b'_1,\dots,b'_{n-1},h(b_n)$ is a
regular sequence mod $h(J)$.
\end{lemma}

\begin{proof}
Induction on $n$.  For $n=1$, $\Gr(b_1)\geq 1$ means that $b_1$ is regular
mod $J$; take $h=\mathrm{id}$.  For $n+1$, in $R[u]$ let
\[
f=b_1+u\,(b_2+u\,b_3+\dots+u^{n-1}b_{n+1}),
\]
the Kronecker polynomial of $(b_1,\dots,b_{n+1})$ with constant term $b_1$.
Since $(b)$ is regular mod $J$ ($\Gr\geq1$), $f$ is regular mod $J$
(McCoy, ctete Theorem 4), and $f\in(b)$.  After the base change to $R[u]$,
Theorem 9 gives $\Gr(b)\geq n$ mod $(J,f)$.  As $b_1=f-u(\dots)$ is in
$(J,f,b_2,\dots,b_{n+1})$, $\Gr(b_2,\dots,b_{n+1})\geq n$ mod $(J,f)$.  By
induction, over $R[u]$ and mod $(J,f)$, a good map $h'$ and operations
give $b'_2,\dots,b'_n,h'(b_{n+1})$ regular mod $(J,f)$.  Replacing $b_1$
by $f$ is one elementary operation ($f-b_1=\sum_{l\geq2}u^{l-1}b_l$), and
the operations on $b_2,\dots,b_{n+1}$ are operations on the whole list
(with coefficient $0$ at the first place).  So $h=h'\circ(R\to R[u])$ and
$b'_1=h'(f)$ give a regular sequence $b'_1,\dots,b'_n,h(b_{n+1})$
mod $h(J)$.
\end{proof}

The polynomials $b'_i$ are ctete's $b_i+u_ib_{i+1}+u_i^2b_{i+2}+\cdots$;
here they arise one at a time, each as one elementary operation.

\section{The induction step}

Assume part 2 for $n$, for all rings and ideals, and let $a,b\in R^{n+1}$,
$V$, $\Gr(b)\geq n+1$ mod $J$.  By Lemma~\ref{lem:generic} and the
invariance properties we may assume (after a good map, renaming $h(J)$
as $J$) that
\[
b_1,\dots,b_n,b_{n+1}\text{ is a regular sequence mod } J .
\]
Write $K=(J,b_1,\dots,b_n)$; so $b_{n+1}$ is regular mod $K$.

\paragraph{Position of $a$.}  If $x$ kills $(a)$ mod an ideal containing
$J$, it kills $(b)$, since $b\equiv aV$.  Hence $(a)$ is regular mod $J$
(as $(b)$ is) and mod $K$ (as $b_{n+1}$ is).  In $R[u]$ let
\[
a'_{n+1}=a_{n+1}+u\,(a_n+u\,a_{n-1}+\dots+u^{n-1}a_1).
\]
It is regular mod $J$ and mod $K$ (McCoy), and replacing $a_{n+1}$ by
$a'_{n+1}$ is an elementary operation on $a$.  We rename $a'$ as $a$.

\paragraph{The grade mod $a_{n+1}$.}  Replacing the last term of a
regular sequence by an element regular modulo the previous terms gives a
regular sequence: so $b_1,\dots,b_n,a_{n+1}$ is a regular sequence
mod $J$, in the ideal $(b_1,\dots,b_n,a_{n+1})$, which therefore has
$\Gr\geq n+1$ mod $J$.  By Theorem 9 with $a_{n+1}$, it has $\Gr\geq n$
mod $(J,a_{n+1})$, hence, as $a_{n+1}\in(J,a_{n+1})$,
\[
\Gr(b_1,\dots,b_n)\geq n\ \text{mod } J'=(J,a_{n+1}).
\]
Modulo $J'$, $b_j\equiv\sum_{i\leq n}a_iV_{ij}$ for $j\leq n$, that is,
$b_{<}\equiv a_{<}W$ with $W$ the upper left $n\times n$ block of $V$;
let $\delta=\det W$.  By induction, mod $(J',b_<)=(K,a_{n+1})$:
$x\delta\in(K,a_{n+1})$ implies $x\in(K,a_{n+1},a_<)$, and
$ya_i\in(K,a_{n+1})$ for all $i\leq n$ implies $y\in(K,a_{n+1},\delta)$.

\paragraph{Cramer.}  By part 1, $\Delta a_{n+1}\equiv\sum_jb_j\adj(V)_{j,n+1}$
mod $J$, and $\adj(V)_{n+1,n+1}=\delta$; hence
\[
\Delta a_{n+1}\equiv\delta\, b_{n+1}\pmod K .
\]

\paragraph{First inclusion.}  Let $x\Delta\equiv b_{n+1}y$ mod $K$.  Then
$b_{n+1}x\delta=x\delta b_{n+1}\equiv x\Delta a_{n+1}\equiv b_{n+1}a_{n+1}y$
mod $K$, so $x\delta\equiv a_{n+1}y$ mod $K$ as $b_{n+1}$ is regular mod
$K$.  Thus $x\delta\in(K,a_{n+1})$, and by induction
$x\in(K,a_{n+1},a_<)\subseteq(J,b,a)$.

\paragraph{Second inclusion.}  Let $xa_{n+1}\equiv b_{n+1}y$ and
$xa_i\equiv b_{n+1}y_i$ mod $K$ for $i\leq n$.  Then
\[
b_{n+1}ya_i\equiv a_ixa_{n+1}=a_{n+1}xa_i\equiv b_{n+1}a_{n+1}y_i\pmod K,
\]
so $ya_i\equiv a_{n+1}y_i$, i.e.\ $ya_i\in(K,a_{n+1})$, for all $i\leq n$.
By induction $y\equiv a_{n+1}t+\delta s$ mod $K$.  Then
\[
a_{n+1}x\equiv b_{n+1}y\equiv a_{n+1}b_{n+1}t+\delta b_{n+1}s\equiv
a_{n+1}(b_{n+1}t+\Delta s)\pmod K,
\]
and since $a_{n+1}$ is regular mod $K$, $x\equiv b_{n+1}t+\Delta s$ mod $K$:
$x\in(J,b,\Delta)$.

\paragraph{Back to the data.}  The conclusions hold for the changed data
in the extended ring; by the invariance properties they hold for the
images of $a,b,\Delta$ (elementary operations keep $(J,b)$, put $a'$ in
$(J,a)$ and keep $\Delta$), and then for $a,b,\Delta$ in $R$ (good map).
This ends the induction and the proof of Theorem~\ref{thm:wiebe}.

\begin{remark}
Compared with ctete: the statement is relative to $J$ and the ring
$T=R/(a_n)$ becomes the ideal $(J,a_{n+1})$, which is what makes the
induction go through without quotient rings; the claim
$\Gr(b_1,\dots,b_n)\geq n$ mod $(J,a_{n+1})$ is derived explicitly by
replacing the last term of the generic regular sequence by $a'_{n+1}$
before applying Theorem 9; and ``without changing $\Delta$'' is the fact
that each elementary matrix has determinant~$1$.
\end{remark}

\section{The formalisation}

The proof is checked in \code{sag/wiebe} (168 modules with the linked
foundation, about 2.6\,s; each of the eight files below checks in under
0.1\,s, and no declaration opens more than 226 definitions).  It uses the
true grade and Theorem 9 of \code{sag/koszul}, McCoy modulo an ideal from
\code{sag/cffr}, and determinants, the adjugate, transposition and the
naturality of $\det$ from \code{sag/exterior}.  Nothing is assumed.

Matrices are lists of columns, and the hypothesis is
$b_j\equiv\sum_ia_iV\,i\,j$; the columns of the formal $V$ are the rows of
the $V$ above, which makes Cramer's rule the identity
$\adj(V)V=\det V\cdot I$ without transposition.  The lists $(J,c)$ are
built with the last generator outermost, so that $x\in(L,c)$ iff
$x-cw\in L$ for some $w$.

\medskip
\begin{tabular}{p{0.32\textwidth}p{0.6\textwidth}}
\code{WiebeLists} & the lists $(J,c)$, inclusion by generators, congruence mod an ideal \\
\code{WiebeMatrices} & elementary operations and $\det E=1$; Cramer at the last index; $\adj(V)_{nn}=\det W$; the block hypothesis \\
\code{WiebeCore} & the two inclusions at $n+1$ from $n$ (\code{Step}), and $n=0$ \\
\code{WiebeSequences} & last term of a regular sequence, replacing it, changing the ideal; good maps \\
\code{WiebeChain} & chains of elementary operations: ideals, $V$ and $\det$, along ring maps \\
\code{WiebeGeneric} & Lemma~\ref{lem:generic} (\code{generic}) \\
\code{WiebeTransfer} & invariance of the conclusions; descent along good maps \\
\code{Wiebe} & the induction (\code{theorem11}); the equalities (\code{Thm11.WithGrade}); Corollary 12 (\code{Thm11.WithSequence})
\end{tabular}
\medskip

Lemma~\ref{lem:generic} quantifies over the extension ring; it is stated
as ``for every proposition $P$, if every such extension gives $P$, then
$P$'', so that no type of rings has to be truncated.

\end{document}
